\documentclass[aspectratio=169]{beamer}  % artikeltyp, legt vorhandene Ebenen (chapter, section etc.)

% =Sprache=
%\usepackage[USenglish,british,american,australian,english]{babel}   %Für Englisch
\usepackage[ngerman]{babel}   % Für deutsch
\usepackage[utf8]{inputenc}   % Gibt utf8 also umlaute, ß etc
\usepackage{lmodern} % update im der Latex schrift, macht schriftbild etwas schöner
\usepackage[T1]{fontenc}   % läd zeichensätze für westeuropäische sprachen nicht unbedingt nötig

% =Mathe=
\usepackage{amsmath} % Grundlegendes packet, gibt Matheumgebungen
\usepackage{amssymb} % Gibt Mathesymbole (pfeile etc) http://milde.users.sourceforge.net/LUCR/Math/mathpackages/amssymb-symbols.pdf
\usepackage{dsfont}  % double struck fonts (z.B: für Mengen von Zahlen) http://milde.users.sourceforge.net/LUCR/Math/mathpackages/dsfont-symbols.pdf
\usepackage{mathrsfs}   % Geschwungene Buchstaben http://www.stat.colostate.edu/~vollmer/pdfs/typesetting-script.pdf
% es gibt da noch seeehr viele mehr...
\usepackage{nicefrac}
\usepackage{siunitx}
\usepackage{mathdots}

% =Bilder=
\usepackage{graphicx}   % Bilder einbinden
\usepackage{multirow}
%\usepackage{caption}

\usepackage{pdfpages}
\usepackage{amsmath}

% =Malen=
%\usepackage[colorlinks]{hyperref}
\usepackage{tikz} % Tool zum zeichnen von dingen
\usepackage{pgf-pie}
\usepackage{pgfplots, pgfplotstable}   % functions


% Pakete aus tikz, die geladen werden (gibt noch viel mehr).
\usetikzlibrary{positioning}
\usetikzlibrary{arrows,automata,shapes}
\usetikzlibrary{trees,fit,decorations.pathreplacing}  % group
\usetikzlibrary{calc} % calculate for relative positioning
\usetikzlibrary{graphs} %  graph placement
\usetikzlibrary{matrix,backgrounds} % Benötigt für hintergründe bei matrizen über tikz
% http://tex.stackexchange.com/questions/57152/how-to-draw-graphs-in-latex

\tikzset{
  invisible/.style={opacity=0},
  visible on/.style={alt={#1{}{invisible}}},
  alt/.code args={<#1>#2#3}{%
    \alt<#1>{\pgfkeysalso{#2}}{\pgfkeysalso{#3}}
  },
}


% =Optimierungen=
\graphicspath{{./Bilder/}} % Gibt den Standard ordner für bilder an
\usepackage{enumerate}  % gives extra options for \begin{enumerate}, e.g. use a),b) instead of 1),2)
\usepackage{hyperref}   % Klickbare Links
\usepackage{caption} % Gives options for customising captions
\usepackage{subcaption} % different captions for multiple pictures in one float
%\usepackage{placeins}  % gibt \FloatBarrier. Wenn das benötigt wird, leben überdenken.
\usepackage{float,rotating}   % Option for fix possitions; rotating of Graphics/tables
\usepackage{array,booktabs}   % erweiterte Möglichkeiten zur Tabellendarstellung
\usepackage{longtable}  % gibt optionen für tabellen über mehrere Seiten
%\usepackage{uline}  % allowes underlines, its strongly recomendet not to underline
%\setlength{\parindent}{0cm}  % Einrücktiefe auf 0cm setzen (siehe weiter unten im Dokument)
\usepackage{indentfirst}
\usepackage[justification=centering]{caption}

\newcommand{\MYhref}[3][blue]{\href{#2}{\color{#1}{#3}}}%
\usepackage{titlesec}

\usepackage[utf8]{inputenc}
\usepackage[T1]{fontenc}
\usepackage{textcomp}
\usepackage{gensymb}

\usepackage{stmaryrd}

\usepackage{rotating}

\usepackage{multicol}

\usepackage[linewidth=1pt]{mdframed}


\usepackage{makeidx}  % gives an Index
\makeindex   % Starts looking for index-keys
\newcommand{\Index}[1]{#1\index{#1}} % gives the comand \Indext{test} which puts "`test"' in the Index, refering to the page the key is on.

% =Definitionen= (Siehe unten im Matheteil)
\newcommand{\pd}{\partial}
\newcommand{\pfrac}[2]{\frac{\partial #1}{\partial #2}}
%%%
% BODY
%%%
\setlength{\parindent}{0em}

\usepackage{listings}

\lstset{ % add your own preferences
    frame=none,
    basicstyle=\footnotesize,
    keywordstyle=\color{Blue1},
    numbers=none,
    showstringspaces=false, 
    stringstyle=\color{Blue3},
    tabsize=2,
    % xleftmargin=2em,
    escapeinside=``,
    moredelim=**[is][\only<+>{\color{black}\lstset{style=highlight}}]{@}{@},
    language=Go, % this is it !
    morekeywords={function, of, number}
}

\usepackage{amsthm}

\usepackage{tikz}

\usepackage{aligned-overset}

\usepackage[ruled,vlined,noend]{algorithm2e}

\usepackage[backend=biber, style=ieee]{biblatex}
\addbibresource{references.bib}
\usepackage{csquotes}

%%%
% notes
%%%
\usepackage{pgfpages}
% \setbeameroption{show only notes}
% \setbeameroption{show notes}
% \setbeameroption{show notes on second screen}



% zu Zeigen
\newcommand{\zz}{$\mathrm{\raise0.25ex\hbox{Z}\kern-.56em\raise-0.25ex\hbox{Z}}$ }

\definecolor{Blue1}{RGB}{0, 74, 153}
\definecolor{Blue2}{RGB}{42, 110, 187}
\definecolor{Blue3}{RGB}{111, 154, 211}
\definecolor{Red1}{RGB}{193, 0, 42}
\definecolor{Red2}{RGB}{123, 41, 39}
\definecolor{GreyDark}{RGB}{54, 53, 52}
\definecolor{GreyLight}{RGB}{242, 243, 241}

%customize itemize
\mode<presentation>
{
    \usetheme{Madrid}
    \usecolortheme{dolphin}

    \setbeamercolor{title}{parent=structure,fg=GreyLight,bg=Blue1}
    \setbeamercolor*{palette primary}{use=structure,fg=GreyDark,bg=Blue3}
    \setbeamercolor*{palette secondary}{use=structure,fg=GreyLight,bg=Blue2}
    \setbeamercolor*{palette tertiary}{use=structure,fg=GreyLight,bg=Blue1}
    \setbeamercolor{frametitle}{parent=structure,bg=GreyLight,fg=Blue1}

    \usefonttheme{professionalfonts}
    \setbeamercolor{itemize item}{fg=Blue1}
    \setbeamertemplate{itemize item}[circle]
    \setbeamercolor{itemize subitem}{fg=Blue2}
    \setbeamertemplate{itemize subitem}[circle] 
    \setbeamercolor{itemize subsubitem}{fg=Blue3}
    \setbeamertemplate{itemize subsubitem}[circle]
    
    \beamertemplatenavigationsymbolsempty
}

\author[Erik Kassubek]{Erik Kassubek} 
\title[Dynamic Analysis of Concurrent Go-Programs]{Dynamic Analysis of Concurrent Go-Programs}
\subtitle{Bachelorarbeit}
\institute[]{Institut für Informatik\\Albert-Ludwigs-Universität Freiburg}
\date{15.02.2023}

\begin{document}

\frame{
    \titlepage
}



% \begin{frame}{Nebenläufigkeit in Go}
%     \begin{itemize}
%       \item Go-Routine
%       \begin{itemize}
%         \item leichtgewichtiger Thread
%         \item ermöglicht Nebenläufigkeit
%       \end{itemize}
%       \item Mutexe (Locks)
%       \begin{itemize}
%         \item Synchronisationsmechanismus
%         \item Löst das Problem des gegenseitigen Ausschluss
%       \end{itemize}
%       \item Channel
%       \begin{itemize}
%         \item Synchronisationsmechanismus
%         \item Ermöglichen Kommunikation zwischen Routinen
%       \end{itemize}
%     \end{itemize}
%     \note{Note}
% \end{frame}

% ============== Mutex ===============

\begin{frame}{Mutex}
  \begin{minipage}{0.49\textwidth}
    \lstinputlisting{code/mutex-example.txt}
  \end{minipage}
  \hfill
  \uncover<2>{\vrule{}}
  \hfill
  \begin{minipage}{0.49\textwidth}
    \uncover<2>{\lstinputlisting{code/mutex-rw.txt}}
  \end{minipage}
\end{frame}

\begin{frame}{Mutex}
  \begin{minipage}[t]{0.49\textwidth}
    \lstinputlisting{code/mutex-zyclic.txt}
  \end{minipage}
  \hfill
  \uncover<2>{\vrule{}}
  \hfill
  \begin{minipage}[t]{0.49\textwidth}
    \uncover<2>{\lstinputlisting{code/mutex-double.txt}}
  \end{minipage}
\end{frame}


% ============== Channel =============

\begin{frame}{Channel - Unbuffered}
  \begin{minipage}{0.49\textwidth}
    \lstinputlisting{code/channel-un-example.txt}
  \end{minipage}
  \hfill
  \uncover<2>{\vrule{}}
  \hfill
  \begin{minipage}{0.49\textwidth}
    \uncover<2>{\lstinputlisting{code/channel-un-block.txt}}
  \end{minipage}
\end{frame}

\begin{frame}{Channel - Buffered}
  \begin{minipage}{0.49\textwidth}
    \lstinputlisting{code/channel-buf-example.txt}
  \end{minipage}
  \hfill
  \uncover<2>{\vrule{}}
  \hfill
  \begin{minipage}{0.49\textwidth}
    \uncover<2>{\lstinputlisting{code/channel-buf-block.txt}}
  \end{minipage}
\end{frame}

\begin{frame}{Channel - Close}
  \begin{minipage}{0.49\textwidth}
    \begin{itemize}
      \item Schließt Channel $\Rightarrow$ keine weiter Kommunikation möglich
      \item Send auf geschlossenem Channel $\Rightarrow$ Laufzeitfehler
    \end{itemize}
  \end{minipage}
  \hfill
  \vrule{}
  \hfill
  \begin{minipage}{0.49\textwidth}
    \lstinputlisting{code/channel-ex-close.txt}
  \end{minipage}
\end{frame}

\begin{frame}{Channel - Select}
  \begin{minipage}{0.49\textwidth}
    \begin{itemize}
      \item Wartet gleichzeitig auf mehrere Channel-Operationen
      \item Erste ausführbare Operation wird ausgeführt
      \item Wenn mehrere gleichzeitig $\Rightarrow$ zufälliger Case
    \end{itemize}
  \end{minipage}
  \hfill
  \vrule{}
  \hfill
  \begin{minipage}{0.49\textwidth}
    \lstinputlisting{code/select-example.txt}
  \end{minipage}
\end{frame}


% ========== Main =============
\begin{frame}{Analyse}
  \begin{itemize}
    \item Entwicklung und Implementierung eines Detektors zur
    \begin{itemize}
      \item Erkennung von problematischen Situationen
      \item Analyse von problematischen Situationen 
    \end{itemize}
    \item<2> Dynamische Analyse
    \begin{itemize}
      \item<2> Instrumentierung
      \item<2> Mehrfache Ausführung für verschiedene Select-Pfade
      \item<2> Verhalten aufzeichnen $\Rightarrow$ Trace
      \item<2> Analyse des Trace nach jeder Ausführung
      \item<2> Konsolidierung der Ergebnisse am Ende
      \item<2> Automatischer Abbruch wenn Laufzeit zu lange
    \end{itemize}
  \end{itemize}  
\end{frame}

% ======== Trace ============ %

% \begin{frame}{Trace}
%   \begin{itemize}
%     \item Speichert Ablauf der notwendigen Informationen
%     \item Ein Trace pro Routine
%     \item Ermöglicht Rekonstruktion des Programmablaufs
%   \end{itemize}
% \end{frame}


% ============== Instrumentierung ====================

\begin{frame}{Instrumentierung}
  \begin{itemize}
    \item Veränderung des Programmcodes
    \item Ersetze Mutexe/Channel durch eigene Objekte
    \item Funktionen auf Objekten:
    \begin{itemize}
      \item Ausführung der eigentlichen Operation
      \item Aufzeichnung der Operation
    \end{itemize}
    \item Aufzeichnung von Fork
    \item Kontrolle über Select
  \end{itemize}
\end{frame}



\begin{frame}{Instrumentierung}
  \begin{itemize}
    \item Funktionen in Package implementiert
    \item Automatisierte Ersetzung im Programmcode
    \begin{itemize}
      \item Traversierend des AST
      \item Veränderung des AST wenn notwendig
    \end{itemize}
    \item Erzeugung einer neuen Main-Datei zur Ausführung
  \end{itemize}
\end{frame}



% ================= Analyse =====================

\begin{frame}{Analyse - Mutexe}
  \begin{itemize}
    \item Suche nach doppeltem Locking
    \begin{itemize}
      \item Lock-Operation als letztes Element in Trace
      \item Traversiere Trace rückwärts
      \item Wenn Lock vor Unlock auf selben Mutex $\Rightarrow$ doppeltes Locking
      \item Kein Blocking wenn beide RLock
    \end{itemize}
  \end{itemize}
\end{frame}

\begin{frame}{Analyse - Mutexe}
  \begin{itemize}
    \item Suche nach zyklischem Locking
    \begin{itemize}
      \item Aufbau von Lock-Bäumen
      \item Suchen nach Zyklen in Lock-Bäumen
    \end{itemize}
  \end{itemize}
\end{frame}

\begin{frame}{Detektion von Deadlocks: Lock-Bäume}
  \begin{minipage}{0.65\textwidth}
      \lstinputlisting[multicols=2]{code/lock-tree.txt}
  \end{minipage}
  \hfill
  \vrule{}
  \hfill
  \begin{minipage}{0.25\textwidth}
      \begin{tikzpicture}[node distance={12mm}, thick, main/.style = {draw}] 
          \node[main, rectangle] (R1) {R1};
          \node[main, circle] [below of=R1](y1) {y}; 
          \node[main, circle] [below of=y1](z1) {z};
          \node[main, circle] [right of=z1](x1) {x};
          \node[main, rectangle] [below of=z1](R2) {R2};
          \node[main, circle] [below of=R2](z2) {x};
          \node[main, circle] [below of=z2](x2) {z};
          \node[main, rectangle] [below of=x1](R3) {R3};
          \node[main, circle] [below of=R3](x3) {z};
          \node[main, circle] [below of=x3](y3) {y};
          
          \draw[->] (R1) to (y1);
          \draw[->] (y1) to (x1);
          \draw[->] (y1) to (z1);
          \draw[->] (R2) to (z2);
          \draw[->] (z2) to (x2);
          \draw[->] (R3) to (x3);
          \draw[->] (x3) to (y3);

          \draw[->, draw=Red1, visible on=<2>] (x3) to (y3);
          \draw[->, draw=Red1, visible on=<2>] (y1) to (x1);
          \draw[->, draw=Red1, visible on=<2>] (z2) to (x2);


          \draw[dashed, draw=Red1, ->, visible on=<2>] (x1) to [out=180+45,in=45] (z2);
          \draw[dashed, draw=Red1, ->, visible on=<2>] (x2) to [out=45,in=180+45] (x3);
          \draw[dashed, draw=Red1, ->, visible on=<2>] (y3) to [out=45,in=0,looseness=1.4] (y1);
      \end{tikzpicture}
  \end{minipage}
\end{frame}

\begin{frame}{Analyse - Mutex}
  \begin{itemize}
    \item Zyklus nur gültig wenn:
    \begin{itemize}
      \item keine zwei R-Lock hintereinander
      \item Locks in Zyklus dürfen keine gemeinsamen Vorfahren haben, außer beide 
        Vorfahren sind RLock
    \end{itemize}
  \end{itemize}
\end{frame}

\begin{frame}{Analyse - Channel}
  \begin{itemize}
    \item Suche nach potenziellem Ablauf mit 
    \begin{itemize}
      \item Send/Receive ohne Kommunikationspartner
      \item potenzielles Send auf geschlossenem Channel
    \end{itemize} 
  \end{itemize}
\end{frame}

\begin{frame}{Analyse - Channel}
  \begin{itemize}
    \item Berechne Pre- und Post-Vectorclocks
    \item Bestimme vectorclock-annotated Trace
    \item Bestimme potenzielle Kommunikationspartner
    \item Betrachte potenzielle Abläufe 
    \item Suche nach Send/Receive ohne Kommunikationspartner und Send auf Closed
  \end{itemize}
\end{frame}

\begin{frame}{Analyse - Channel}
  \lstinputlisting{code/vc-code.txt}  
\end{frame}

% \begin{frame}{Analyse - Channel}
%   \begin{align*}
%     [&[signal(1, 2), signal(2, 3), pre(3, 1?), post(9, 1, 1?, 6)]\\
%     &[wait(5, 2), pre(6, 1!), post(7, 1, 1!), pre(8, 1?), post(12, 2, 1?, 10)]\\
%     &[wait(4, 3), pre(10, 1!), post(11, 2, 1!)]]
%   \end{align*}
% \end{frame}

% \begin{frame}{Analyse - Channel}
%   \begin{minipage}{0.49\textwidth}
%   \centering
%   \includegraphics[width=0.6\textwidth]{img/vc2.eps}
%   \end{minipage}
% \end{frame}


\begin{frame}{Analyse - Channel}
  \begin{minipage}{0.49\textwidth}
  \centering
  \includegraphics[width=0.6\textwidth]{img/vc.eps}
  \end{minipage}
  \hfill
  \uncover<2->{\vrule{}}
  \hfill
  \begin{minipage}{0.49\textwidth}
    \uncover<2->{
      \begin{align*}
        [&[^{[2,0,0]}x?^{[3,2,0]}]\\
        &[^{[1, 1, 0]}x!^{[1, 2, 0]}, ^{[1, 2, 0]}x?^{[2, 3, 2]}]\\
        &[^{[2, 0, 1]}x!^{[2, 0, 2]}]]
      \end{align*}
    }
    \begin{itemize}
      \item<3> Für Operation ohne Post:\\Post-VC: [max, ..., max]
    \end{itemize}
  \end{minipage}
\end{frame}

\begin{frame}{Analyse - Channel}
  \begin{itemize}
    \item $x$ ist Ursache von $y$ wenn
    \begin{itemize}
      \item $x \to y \Leftrightarrow VC(x) < VC(y) \Leftrightarrow \forall z (VC(x)_z \leq VC(y)_z) \land \exists z' (VC(x)_{z'} < VC(y)_{z'})$  
    \end{itemize}
    \item $x$ ist nebenläufig/unvergleichbar mit $y$ wenn
    \begin{itemize}
      \item $ x || y \Leftrightarrow VC(x) \not \lessgtr VC(y) \Leftrightarrow VC(x) \nless VC(y) \land VC(y) \nless VC(x)$
    \end{itemize}
  \end{itemize}
\end{frame}

\begin{frame}{Analyse - Channel}
  \begin{itemize}
    \item Potenzielle Kommunikationspartner auf ungebufferten Channels:
    \begin{itemize}
      \item Paar von Send-Receive auf selbem Channel
      \item Pre- oder Post-Vectorclock unvergleichbar
      \item Keine gemeinsamen gehaltenen Mutexe bei Ausführung
    \end{itemize}
    \uncover<2>{\item Potenzielle Kommunikationspartner auf gebufferten Channels:
    \begin{itemize}
      \item Paar von Send-Receive auf selbem Channel
      \item $\#s_{<} \leq \#r_{<} + \#r_{\not\lessgtr}$
      \item $\#s_{<} + \#s_{\not\lessgtr} \geq \#r_{<}$
    \end{itemize}
    }
  \end{itemize}
\end{frame}

\begin{frame}{Analyse - Channel}
  \lstinputlisting{code/buffered_form.txt}
\end{frame}


\begin{frame}{Analyse - Channel}
  \begin{minipage}{0.49\textwidth}
    \begin{align*}
      [&[^{[2,0,0]}x_1?^{[3,2,0]}]\\
      &[^{[1, 1, 0]}x_2!^{[1, 2, 0]}, ^{[1, 2, 0]}x_3?^{[2, 3, 2]}]\\
      &[^{[2, 0, 1]}x_4!^{[2, 0, 2]}]]
    \end{align*}
    \uncover<2->{
      \begin{align*}
        &x_2! \to x_1?\\
        &x_4! \to x_1?\\
        &x_4! \to x_3?
      \end{align*}
    }
    \end{minipage}
    \hfill
    \uncover<3->{\vrule{}}
    \hfill
    \begin{minipage}{0.24\textwidth}
      \uncover<3->{
        \begin{align*}
          &x_2! \to x_1?\\
          &x_4! \to x_3?
        \end{align*}  
      }
      \uncover<4->{
        \begin{align*}
          &x_4! \to x_1?\\
          &x_2! \to \lightning
        \end{align*}  
      }
    \end{minipage}
    \hfill
    \begin{minipage}{0.24\textwidth}
      \uncover<5->{
        \begin{align*}
          &x_1? \leftarrow x_2!\\
          &x_3? \leftarrow x_4!
        \end{align*}  
      }
      \uncover<6->{
        \begin{align*}
          &x_1? \leftarrow x_4!\\
          &x_3? \leftarrow \lightning
        \end{align*}  
      }
    \end{minipage}
\end{frame}

\begin{frame}{Analyse - Channel}
  \begin{itemize}
    \item Potenzielles Send auf Closed Channel:
    \begin{itemize}
      \item Pre- oder Post-VC von Send unvergleichbar mit VC von Close
      \item Wenn Close ist Ursache von Send $\Rightarrow$ Send auf Close tritt 
      während Durchlauf auf $\Rightarrow$ erzeugt Laufzeitfehler
    \end{itemize}
  \end{itemize}
\end{frame}

\begin{frame}{Zusammenfassung}
  \begin{itemize}
    \item Dynamischer Detektor zur Erkennung von Concurrency-Bugs
    \begin{itemize}
      \item Mutex-Bugs mit Lock-Bäumen
      \item Channel-Bugs mit Vectorclock
    \end{itemize}
    \item Implementierung für Instrumenter und Detektor:
    \begin{itemize}
      \item \url{https://github.com/ErikKassubek/GoChan}
    \end{itemize}
  \end{itemize}
\end{frame}





% ===== Quellen =====

\begin{frame}{Sources}
    \nocite{*}
    \AtNextBibliography{\tiny}
    \printbibliography[heading=none, keyword={1}]
\end{frame}

\begin{frame}{Sources}
  \nocite{*}
  \AtNextBibliography{\tiny}
  \printbibliography[heading=none, keyword={2}]
\end{frame}



% ========== Appendix ==========

\begin{frame}[noframenumbering]{Instrumentierung - Fork}
    \begin{itemize}
      \item \texttt{signal(ts, nr)} in Trace
      \item Erzeuge neuen Trace
      \item Starte der neuen Routine
      \item \texttt{wait(ts, nr)} in neuen Trace
      \item Ausführen der Routine
    \end{itemize}
\end{frame}

\begin{frame}[noframenumbering]{Instrumentierung - Mutex}
  \begin{itemize}
    \item (Try)(R)Lock:
    \begin{itemize}
      \item \texttt{lock(ts, mu.id, ⟨-|t|r|tr⟩, ⟨0|1⟩)}
      \item Führe Lock aus
    \end{itemize}
    \item (Try)(R)Unlock:
    \begin{itemize}
      \item \texttt{unlock(ts, mu.id)} in Trace
      \item Führe Lock aus
    \end{itemize}
  \end{itemize}
\end{frame}

\begin{frame}[noframenumbering]{Instrumentierung - Channel}
  \begin{itemize}
    \item Send: (c <- i):
    \begin{itemize}
      \item \texttt{pre(ts, c.id, true)} in Trace
      \item \texttt{c <- \{i, r.id, ts\}}
      \item \texttt{post(ts', c.id, true, ts)} in Trace
    \end{itemize}
    \item Receive (a := <- c):
    \begin{itemize}
      \item \texttt{pre(ts, c.id, false)} in Trace
      \item \texttt{\{i, id\_send, ts\_send\} := <-c}
      \item \texttt{post(ts', ch.id, false, id\_send, ts\_send)} in Trace
      \item \texttt{a := i}
    \end{itemize}
    \item Close (close(c))
    \begin{itemize}
      \item \texttt{close(ts, c.id)} in Trace
      \item \texttt{close(c)}
    \end{itemize}
  \end{itemize}
\end{frame}




\begin{frame}[noframenumbering]{Instrumentierung - Select}
  \begin{tabular}{lcl}
    select(e$_\text{i} \leadsto \text{B}_\text{i}$) & $\Rightarrow$ & ts := atomicInc(c); [ pre(ts, e$_1^*$, $\ldots$, e$_n^*$, false) ]$^\text{id}$;\\
    & & select(e$_\text{i} \leadsto$ \{ ts' := atomicInc(c);\\
    & & [ $\langle \text{post(ts, e$_\text{i}$.ch, false, e$_\text{i}^\text{m}$.id\_send, e$_\text{i}^\text{m}$.ts\_send)}\ |$ \\
    & & post(ts, e$_\text{i}$.ch, true, id) $\rangle$ ]$^\text{id}$ B$_\text{i}$\}) \\
  select(e$_\text{i} \leadsto \text{B}_\text{i}$ | B$_\text{def}$) & $\Rightarrow$ & ts := atomicInc(c); [ pre(ts, e$_1^*$, $\ldots$, e$_n^*$, false) ]$^\text{id}$;\\
    & & select(e$_\text{i} \leadsto$ \{ ts' := atomicInc(c);\\
    & & [ $\langle \text{post(ts, e$_\text{i}$.ch, false, e$_\text{i}^\text{m}$.id\_send, e$_\text{i}^\text{m}$.ts\_send)}\ |$ \\
    & & post(ts, e$_\text{i}$.ch, true, id) $\rangle$ ]$^\text{id}$ B$_\text{i}$\} |\\
    & & ts' := atomicInc(c); [ default(ts) ]$^\text{id}$; B$_\text{def}$) 
  \end{tabular}
\end{frame}

\begin{frame}[noframenumbering]{Instrumentierung - Select}
  \lstinputlisting{code/select-pre.txt}
\end{frame}

\begin{frame}[noframenumbering]{Instrumentierung - Select}
  \lstinputlisting{code/select-post.txt}
\end{frame}


\end{document}
